%% kapa: exam=1 round=34 year=2023 period=1 form=A subject=real_estate subject_ko=부동산학원론 seq=34 number=114
%% source: Q-Net 감정평가사 기출문제 (cst00309) artlSeq=5211789 file=raw/1008/5211789/2023년도 제34회 감정평가사 1차시험 1교시.pdf
%% conversion: kordoc 4.24.51 -> tools/gen1.py (auto); review=proofread-sonnet; flags=none
\begin{question}{114}
토지에 관한 강제경매절차에서 토지의 부합물로서 낙찰자가 소유권을 취득할 수 있는 경우를 모두 고른 것은? (다툼이 있으면 판례에 의함)
\begin{framedtext}
ㄱ. 토지소유자가 마당에 설치한 연못 \par
ㄴ. 타인이 토지소유자의 동의 없이 임의로 심은 조경수 \par
ㄷ. 토지에 지상권을 가진 자가 경작을 위해 심은 감나무 \par
ㄹ. 기둥, 지붕 및 주벽의 공사가 완료되어 건물로서의 외관을 갖추었으나 사용승인을 받지 못한 건물 \par
\end{framedtext}
\begin{choices}
  \item ㄱ, ㄴ
  \item ㄴ, ㄷ
  \item ㄱ, ㄴ, ㄷ
  \item ㄱ, ㄷ, ㄹ
  \item ㄱ, ㄴ, ㄷ, ㄹ
\end{choices}
\end{question}
\answer{1}
